  \documentclass[12pt]{article}

  \usepackage{sbc-template}
  \usepackage{graphicx,url}
  \usepackage[utf8]{inputenc}
  \usepackage[brazil]{babel}
  \usepackage{xcolor}

  \sloppy

  \title{Pipeline Baseada em LLM para Análise Automatizada de Editais de Fomento}

  \author{Gabriel Lopes de Albuquerque, Augusto César Ferreira de Miranda Oliveira,\\ Francisco de Assis Carlos Filho}

  \address{Universidade de Pernambuco (UPE) -- Campus Surubim\\
    Surubim -- PE -- Brasil
    \email{gabriel.lalbuquerque@upe.br, augusto.oliveira@upe.br, francisco.assis@upe.br}
  }

  \begin{document}

  \maketitle

  \begin{abstract}
    Manually monitoring public funding calls, spread across portals with different formats, is costly and makes it hard for each call to reach the researchers in its area in time. This work proposes a pipeline that collects calls from FACEPE, CNPq, CAPES, and FINEP, extracts the text of the PDFs (using OCR when needed), and uses the Gemini 2.5-Flash LLM, guided by a structured prompt, to produce data that route each call to the interested faculty. We present the architecture, the coverage achieved, and the evaluation protocol for the hypotheses of automatic targeting (H1) and extraction reliability (H2).
  \end{abstract}

  \begin{resumo}
    O acompanhamento manual de editais de fomento público, distribuídos em portais com formatos distintos, é custoso e dificulta que cada edital chegue a tempo aos pesquisadores de sua área. Este trabalho propõe uma pipeline que coleta editais da FACEPE, CNPq, CAPES e FINEP, extrai o texto dos PDFs (com OCR quando necessário) e usa o LLM Gemini 2.5-Flash, orientado por prompt estruturado, para gerar dados que direcionam cada edital aos docentes interessados. São apresentados a arquitetura, a cobertura alcançada e o protocolo de avaliação das hipóteses de direcionamento (H1) e de confiabilidade da extração (H2).
  \end{resumo}

  \noindent\textbf{Palavras-chave:} web scraping, extração de informação de documentos, large language models, prompt engineering, data pipeline, editais públicos

  \noindent\textbf{Keywords:} web scraping, document information extraction, large language models, prompt engineering, data pipeline, public funding calls

  \section{Introdução}\label{sec:introducao}

  Instituições de pesquisa e pesquisadores individuais dependem do acompanhamento contínuo de editais de fomento público, isto é, chamadas de financiamento, bolsas e programas de apoio, publicados por agências como a Fundação de Amparo à Ciência e Tecnologia do Estado de Pernambuco (FACEPE), o Conselho Nacional de Desenvolvimento Científico e Tecnológico (CNPq), a Coordenação de Aperfeiçoamento de Pessoal de Nível Superior (CAPES) e a Financiadora de Estudos e Projetos (FINEP). Cada uma dessas instituições publica seus editais em portais próprios, com estruturas de página, frequência de atualização e formatos de documento distintos entre si: em geral, arquivos PDF sem estrutura semântica clara \cite{b4} e, com frequência, sem disponibilização via API, situação semelhante à descrita por Krc e Schlaack [2023] para fontes governamentais que publicam seu conteúdo em páginas web sem oferecer os documentos por meio de API ou de outras tecnologias de web semântica.

  Nesse cenário, o acompanhamento manual de editais distribuídos em múltiplos portais institucionais é uma tarefa repetitiva, sujeita a atrasos e a erros de leitura. A literatura registra esse custo em contextos próximos: a extração manual de dados de documentos é lenta, custosa e sujeita a erros \cite{b4}, pode consumir várias horas de trabalho por dia \cite{b6} e, no monitoramento contínuo de fontes governamentais, chegou a exigir a dedicação integral de um profissional \cite{b3}. No próprio domínio de editais, Holanda \emph{et al.} [2026] relatam que uma universidade pública brasileira ainda reescreve manualmente seus editais, com esforço repetitivo e maior tempo de publicação. O problema é agravado pelo fato de que parte desses PDFs sequer possui camada de texto extraível (documentos escaneados ou gerados por impressão para PDF a partir de visualizadores que bloqueiam a seleção de texto), caso em que a qualidade do reconhecimento óptico de caracteres (OCR) passa a comprometer a precisão da extração \cite{b9}. O resultado é o risco concreto de perda de prazos de submissão e de decisões de elegibilidade tomadas a partir de leitura incompleta dos critérios do edital. Ao mesmo tempo, os modelos de linguagem de grande porte (LLMs) passaram a executar tarefas de extração de informação sem treinamento específico para cada tarefa \cite{b4,b7}, e métodos que integram OCR e LLMs mostram-se mais robustos ao ruído do reconhecimento de texto do que abordagens tradicionais \cite{b6}. É nesse contexto que se insere o desenvolvimento da pipeline proposta, que automatiza a coleta, extração e estruturação dessas informações.

  Diante desse problema, este trabalho investiga duas perguntas de pesquisa. A primeira (PP1) é se é possível, a partir dos dados estruturados extraídos pela pipeline (áreas de interesse, segmentos e público-alvo), direcionar automaticamente cada edital aos pesquisadores, departamentos ou grupos de pesquisa mais aderentes de sua área, de modo que esses interessados possam se organizar com antecedência e a instituição venha a captar mais recursos e ganhar mais visibilidade junto às agências de fomento. A segunda (PP2) é se é possível construir uma pipeline automatizada, baseada em um LLM, capaz de extrair e estruturar com confiabilidade suficiente os campos essenciais de um edital de fomento público, incluindo documentos sem camada de texto extraível, de modo a viabilizar, como pré-condição técnica, o direcionamento automático investigado na PP1, reduzindo ao mesmo tempo o esforço manual de acompanhamento.

  A relevância deste trabalho reside no enfrentamento da assimetria de informação e do esforço manual que historicamente limitam a captação de recursos acadêmicos. Atualmente, o acompanhamento descentralizado e intempestivo de editais resulta na perda de oportunidades e no engajamento reativo de pesquisadores. Para solucionar esse problema, propõe-se uma abordagem de caráter quantitativo articulada em torno de duas hipóteses interdependentes. A Hipótese 1 mensura o impacto final da ferramenta ao validar a eficácia do direcionamento automatizado e personalizado de oportunidades aos grupos de pesquisa e departamentos, e representa o diferencial estratégico e a contribuição principal da pesquisa: a proatividade e a assertividade na gestão de oportunidades institucionais. Por sua vez, a Hipótese 2 estabelece a base metodológica ao avaliar experimentalmente a acurácia da extração automatizada de dados de editais, pré-requisito técnico para a viabilidade do sistema. Dessa forma, a validação quantitativa da sustentação técnica (H2) garante a precisão necessária para que o direcionamento proposto na H1 seja alcançado.

  A Hipótese 1 (H1) é a de que o direcionamento automático de cada edital aos interessados certos de sua área, obtido pelo cruzamento entre os campos de classificação já extraídos pela pipeline (áreas de interesse, segmentos e público-alvo) e um cadastro de perfis institucionais (docentes, departamentos, grupos de pesquisa), permite que esses interessados se organizem com antecedência suficiente para preparar e submeter propostas mais competitivas. Como consequência, espera-se um aumento mensurável no volume de recursos de fomento captados pela instituição e ganho de visibilidade junto às agências financiadoras, em comparação com o cenário de acompanhamento manual e não direcionado.

  A Hipótese 2 (H2) é a de que uma pipeline que combina coleta específica por fonte (web scraping e integração via API), extração de texto com fallback de reconhecimento óptico de caracteres (OCR) e um LLM orientado por prompt estruturado com vocabulário controlado (dado que a qualidade das respostas de um LLM depende diretamente da estrutura e da clareza do prompt \cite{b1,b5}) é capaz de extrair, de forma consistente, os campos essenciais de um edital (público-alvo, critérios de elegibilidade, cronograma e prazo de submissão), atingindo precisão e recall por campo suficientes para sustentar o direcionamento automático da H1 e reduzindo também o esforço manual de monitoramento, sem comprometer significativamente a confiabilidade das informações extraídas. Trata-se de um ganho relevante, porém secundário diante do objetivo de direcionamento e captação de recursos investigado pela H1.

  Assim, o objetivo geral deste trabalho é desenvolver e avaliar uma pipeline automatizada capaz de coletar, extrair e estruturar, por meio de um LLM, informações de editais de fomento público publicados por múltiplas instituições, disponibilizando os resultados de forma consultável e com notificação proativa de prazos, e direcionar automaticamente cada edital aos interessados mais aderentes de sua área a partir dos campos de classificação extraídos (H1; mecanismo implementado na Seção~\ref{sec:direcionamento-automatico-de-editais}), deixando para trabalhos futuros (Seção~\ref{sec:trabalhos-futuros}) a avaliação quase-experimental dos efeitos desse direcionamento. Para alcançá-lo, foram definidos os seguintes objetivos específicos:
  \begin{itemize}
  \item Implementar coletores (scrapers) específicos para cada fonte de dados monitorada (FACEPE, CNPq, CAPES e FINEP), respeitando as particularidades estruturais de cada portal.
  \item Extrair o conteúdo textual de documentos em PDF, incluindo casos sem camada de texto nativa, por meio de reconhecimento óptico de caracteres (OCR).
  \item Utilizar um LLM orientado por prompt estruturado para transformar o texto extraído em dados estruturados, com vocabulário controlado para os campos de classificação (áreas de interesse e segmentos).
  \item Persistir os dados extraídos em um banco de dados não relacional (MongoDB), com controle de duplicidade por documento.
  \item Notificar automaticamente, por e-mail, a existência de novos editais processados e a proximidade de prazos de submissão.
  \item Disponibilizar uma interface web para que pesquisadores marquem editais de interesse, integrada à mesma base de dados utilizada pela pipeline.
  \item Cruzar os campos de classificação extraídos de cada edital com o cadastro de preferências de docentes, direcionando automaticamente os editais novos aos interessados aderentes de sua área.
  \end{itemize}

  O restante deste artigo está organizado da seguinte forma: a Seção~\ref{sec:fundamentacao-teorica-e-trabalhos-relacionados} apresenta os conceitos fundamentais e discute os trabalhos relacionados à extração de informação de documentos com LLMs e à engenharia de prompt. A Seção~\ref{sec:metodologia} descreve a metodologia adotada, incluindo as fontes de dados, a arquitetura do sistema e o pipeline de processamento. A Seção~\ref{sec:resultados} apresenta os resultados obtidos, incluindo uma análise quantitativa da cobertura alcançada. A Seção~\ref{sec:discussao} discute os resultados à luz da literatura e aponta contribuições e limitações da solução. A Seção~\ref{sec:conclusao} conclui o artigo e aponta direções para trabalhos futuros.

  \section{Fundamentação Teórica e Trabalhos Relacionados}\label{sec:fundamentacao-teorica-e-trabalhos-relacionados}

  \subsection{Conceitos Fundamentais}\label{sec:conceitos-fundamentais}
  Os conceitos apresentados nesta subseção acompanham as etapas do próprio pipeline do sistema proposto: a coleta dos documentos (\emph{web scraping}), a estrutura que organiza essa coleta em um fluxo reprodutível (\emph{data pipeline}), a tarefa de converter o conteúdo desses documentos em dados estruturados (\emph{Document Information Extraction}), a tecnologia empregada para realizar essa tarefa (LLMs) e a técnica usada para controlar seu comportamento (\emph{prompt engineering}).

  \textbf{Web scraping} é a técnica de extração automatizada de dados a partir de páginas web, geralmente por meio da análise de sua estrutura HTML \cite{b1}. É a abordagem adotada para três das quatro fontes do sistema proposto (Seção~\ref{sec:base-de-dados}), cujos portais publicam os editais em páginas HTML e não disponibilizam os documentos por meio de uma API. Essa etapa de coleta é a primeira de um fluxo mais amplo, descrito a seguir como \emph{data pipeline}.

  Uma \textbf{data pipeline} (ou fluxo ETL, sigla de \emph{Extract, Transform, Load}) é uma sequência automatizada de etapas que move dados de uma fonte de origem até um destino final, aplicando transformações intermediárias \cite{b3}. No caso do sistema proposto, essas etapas correspondem à coleta (\emph{extract}), à extração de texto e estruturação via LLM (\emph{transform}) e à persistência em banco de dados (\emph{load}), detalhadas na Seção~\ref{sec:pipeline-de-processamento}. A etapa de transformação, em particular, corresponde à tarefa de extração de informação de documentos, discutida a seguir.

  \textbf{Document Information Extraction} (ou \emph{PDF Data Extraction}) refere-se à tarefa de converter o conteúdo de um documento, tipicamente um PDF, em dados estruturados e semanticamente significativos. Moreira-Filho \emph{et al.} [2025] observam que os PDFs não possuem uma estrutura semântica clara, o que dificulta a identificação de elementos como títulos, tabelas e seções, e tratam a extração estruturada como um problema de reconhecimento de entidades nomeadas (NER) no qual os LLMs se mostram eficientes graças à sua capacidade de compreensão contextual. É justamente essa tecnologia, os LLMs, que se descreve na sequência.

  \textbf{Large Language Models (LLMs)} são modelos de linguagem baseados em técnicas de processamento de linguagem natural, treinados por aprendizado auto-supervisionado em um grande corpus de texto \cite{b7}. Sua escala cresceu de forma acentuada nos últimos anos, de centenas de milhões a mais de um trilhão de parâmetros \cite{b7}, e esses modelos apresentam forte capacidade de generalização, executando tarefas variadas, inclusive a extração de informação a partir de texto, sem treinamento específico para cada tarefa \cite{b7}. Em vez de atualizar os parâmetros do modelo, é possível adaptá-lo a uma tarefa nova apenas por meio do prompt \cite{b5}, o que permite que um LLM siga instruções detalhadas em linguagem natural para estruturar um edital a partir apenas do texto e do prompt fornecidos. Controlar esse comportamento por meio do prompt é, por sua vez, o objeto da técnica descrita a seguir.

  \textbf{Prompt engineering} é definida como o processo de criar e ajustar instruções específicas (o \emph{prompt}) para guiar o comportamento de um LLM e obter respostas mais precisas e consistentes \cite{b5,b1}. Entre as estratégias mais estabelecidas estão o \emph{zero-shot prompting} (nenhum exemplo fornecido), o \emph{few-shot prompting} (poucos exemplos fornecidos no próprio prompt) e o \emph{chain-of-thought (CoT) prompting} (instrução para que o modelo explicite passos intermediários de raciocínio antes da resposta final). O prompt utilizado pelo analisador do sistema proposto (Seção~\ref{sec:transformacao}) combina uma instrução de papel (``atue como um analista de editais''), restrições explícitas de formato de saída (JSON) e um vocabulário controlado para dois dos campos extraídos, numa abordagem próxima da técnica de \emph{role prompting} combinada a \emph{few-shot prompting} descrita por Blázquez-Ochando, Prieto-Gutiérrez e Ovalle-Perandones [2025].

  \subsection{Trabalhos Relacionados}\label{sec:trabalhos-relacionados}

  \subsubsection{Método de Seleção dos Trabalhos}
  Os trabalhos discutidos nesta seção foram reunidos por um mapeamento sistemático da literatura recente em torno dos três eixos temáticos mais próximos do sistema proposto: extração de informação de documentos com LLMs, engenharia de prompt aplicada a tarefas estruturadas e pipelines automatizadas de coleta de documentos. A busca foi conduzida no Google Scholar, pela sua ampla cobertura, e complementada por consulta direta a repositórios e bibliotecas digitais (\emph{arXiv}, IEEE Xplore, Wiley Online Library e periódicos da MDPI) e aos anais do Simpósio Brasileiro de Sistemas de Informação (SBSI), combinando os termos ``web scraping'', ``data pipeline'', ``ETL'', ``document information extraction'', ``PDF data extraction'', ``document understanding'', ``large language models'', ``semantic extraction'', ``document classification'' e ``prompt engineering''. Essa varredura inicial identificou cerca de 20 trabalhos aparentemente aderentes ao escopo. Uma primeira etapa de triagem, por leitura de título, resumo e palavras-chave, manteve os trabalhos publicados entre 2023 e 2025 com aplicação prática ou validação experimental sobre uso de LLMs, extração de dados, automação de documentos ou \emph{prompt engineering}, e excluiu duplicatas, trabalhos de abordagem puramente teórica e os voltados à extração de fontes já estruturadas (como APIs), reduzindo o conjunto a cerca de 10 trabalhos. A segunda etapa, de leitura integral, selecionou 6 trabalhos \cite{b1,b2,b3,b4,b5,b6} como base principal da fundamentação. A esse núcleo somaram-se, posteriormente, um \emph{survey} \cite{b7} usado apenas como referência de enquadramento teórico geral sobre características de LLMs (Seção~\ref{sec:conceitos-fundamentais}) e os dois trabalhos do SBSI \cite{b8,b9}, localizados por busca manual nos anais do evento, totalizando os nove trabalhos aqui discutidos. Sahoo \emph{et al.} [2024] foi localizado pelo próprio web scraper do projeto; os demais, por busca manual.

  \subsubsection{Extração de Dados de Literatura Científica com LLMs e KNIME}
  Moreira-Filho \emph{et al.} [2025] partem da constatação de que o crescimento da literatura científica (mais de 2,5 milhões de artigos publicados por ano) torna a extração manual de dados inviável em escala, e descrevem um \emph{workflow}, construído na plataforma \emph{low-code} KNIME, que integra ferramentas de \emph{parsing} de documentos com LLMs para extrair variáveis de publicações científicas e de PDFs em geral, tratando a tarefa como um problema de reconhecimento de entidades nomeadas (NER) no qual o LLM se mostra mais eficiente que abordagens tradicionais graças à sua capacidade de compreensão contextual. O sistema opera em dois modos: um modo texto, que extrai diretamente texto e tabelas, e um modo imagem, que envia páginas renderizadas a um LLM multimodal. A engenharia de \emph{prompt} é apontada pelos autores como fundamental para garantir precisão e consistência e evitar alucinações, com a temperatura fixada em zero para resultados determinísticos e a saída sempre em JSON. Os autores reportam 81,14\% de acurácia no modo texto para publicações científicas e até 98,54\% no modo imagem para PDFs gerais, com Claude 3.5 Sonnet e GPT-4o obtendo o melhor desempenho e a menor incidência de alucinação entre os modelos testados. É o trabalho mais próximo estruturalmente do sistema proposto entre os consultados, e as decisões de \emph{prompt} descritas (restrição de formato de saída e temperatura zero) são próximas das adotadas no analisador do sistema proposto (Seção~\ref{sec:transformacao}).

  \subsubsection{Engenharia de Prompt Aplicada à Extração de Informação de Documentos}
  Chen \emph{et al.} [2025], em estudo publicado na revista \emph{Electronics}, propõem um método de extração de informação baseado na integração de ferramentas de OCR (Amazon Textract, na etapa de pré-processamento do documento) com um LLM, argumentando que abordagens tradicionais são insuficientes para esse tipo de tarefa: expressões regulares carecem de flexibilidade para lidar com múltiplos formatos de documento, e métodos de reconhecimento de entidades nomeadas (NER) exigem grandes volumes de dados rotulados para treinamento, ao passo que ``métodos baseados em LLM demonstram superioridade ao lidar com ruído de OCR'' \cite{b6}. O prompt estruturado utilizado é otimizado por meio de um \emph{Automatic Prompt Engineer} (APE), que gera e seleciona automaticamente a melhor instrução, e refinado por uma técnica de \emph{Instruction Prompt Calibration} (IPC), que combina avaliação humana e avaliação automática do próprio LLM em uma única métrica ponderada. A saída, um JSON estruturado, passa ainda por uma etapa de pós-processamento com correção baseada em expressões regulares e cálculo de uma pontuação de confiança. Os autores reportam precisão média de 95,5\% e acurácia de 91,5\% na extração de campos-chave, superando modelos de resposta a perguntas (QA) tradicionais usados como comparação. Entre as limitações reportadas pelos próprios autores estão o processamento de documentos multilíngues (que se comporta de forma distinta do cenário monolíngue avaliado) e um limite prático de entrada de 16.385 tokens, ambas relevantes para a discussão de limitações do sistema proposto (Seção~\ref{sec:limitacoes}).

  \subsubsection{Engenharia de Prompt para Web Scraping Bibliográfico}
  Blázquez-Ochando, Prieto-Gutiérrez e Ovalle-Perandones [2025] investigam como a engenharia de prompt pode ser usada para que um LLM (ChatGPT-4o) gere, em uma única interação, um web scraper funcional adaptado a um catálogo bibliográfico, tornando o processo acessível a pesquisadores sem conhecimento avançado de programação. A abordagem combina \emph{role prompting} e \emph{few-shot prompting}, e os autores defendem que um prompt estruturado deve conter contexto, objetivo, restrições, exemplos e etapas detalhadas. O scraper gerado dessa forma coletou mais de 55 mil registros, a uma taxa de 66,76 registros por minuto, evidência de que a abordagem é viável em escala. Assim como no sistema proposto, o LLM é acionado uma única vez por tarefa; o trabalho difere por empregar o LLM para \emph{gerar o código} de scraping, e não para \emph{interpretar o conteúdo} de documentos já coletados.

  \subsubsection{Mapeamento Sistemático de Técnicas de Prompt Engineering}
  Sahoo \emph{et al.} [2024] apresentam um mapeamento sistemático das técnicas de \emph{prompt engineering}, organizadas por finalidade (raciocínio, redução de alucinação, entre outras). O levantamento cobre desde o \emph{zero-shot} e o \emph{few-shot prompting} até variações de \emph{chain-of-thought}, além de estratégias de mitigação de alucinação (\emph{Retrieval-Augmented Generation} e \emph{Chain-of-Verification}), de geração automática de prompts (\emph{Automatic Prompt Engineer}) e de uso de ferramentas externas pelo modelo. Os autores observam que persistem desafios como vieses, imprecisões factuais e falta de interpretabilidade. O trabalho serve de referência para posicionar a estratégia de prompt do sistema proposto dentro do espectro mais amplo de técnicas disponíveis.

  \subsubsection{Pipeline de Metadados em Escala a partir de Web Scraping}
  Krc e Schlaack [2023] relatam a construção de uma pipeline para a biblioteca digital FRASER (Federal Reserve Bank of St. Louis), que automatiza a coleta, a transformação e a ingestão de documentos e metadados de fontes governamentais dos Estados Unidos que, segundo os autores, ``apresentam seu conteúdo em páginas web, mas não disponibilizam os metadados e links de documentos via API ou outras tecnologias de web semântica'' \cite{b3}. O processo, originalmente totalmente manual e não projetado para automação, cresceu a ponto de o monitoramento de alertas e a coleta, descrição e ingestão manuais se tornarem, já em 2021, uma posição de dedicação integral. A solução resultante é um programa em Python executado diariamente que faz o \emph{scraping} de uma página, verifica se os itens capturados já existem no repositório e, em caso negativo, transforma e ingere os metadados e o arquivo e envia uma notificação por e-mail; a arquitetura separa o processamento de dados do gerenciamento de dados, define o que coletar por meio de arquivos JSON (\emph{sitemaps} com seletores) e aciona uma função \emph{serverless} (AWS Lambda) a partir de um \emph{endpoint} HTTP. A deduplicação é feita pela comparação do nome de arquivo construído contra os arquivos já armazenados em um \emph{bucket} AWS S3, e uma etapa manual de conferência e aprovação precede a ingestão final. Os autores reportam que o tempo ativo para adicionar novos registros foi reduzido a poucos cliques e que os alertas de falso positivo foram reduzidos a zero. É o trabalho mais próximo do sistema proposto em termos de escopo (coleta contínua e em escala de documentos governamentais heterogêneos), mas não emprega um LLM em nenhuma etapa: a curadoria final fica a cargo de uma conferência humana.

  \subsubsection{Coleta Multi-Fonte com Ranqueamento Semântico de Relevância}
  Fahrudin \emph{et al.} [2025] partem de um sistema anterior dos próprios autores que dependia de um único site de busca (Google Scholar, via automação com Selenium) e de execução em thread única, o que, segundo os autores, ``frequentemente encontra um número limitado de artigos de referência relevantes após um longo tempo'' \cite{b2}. O sistema proposto introduz três melhorias sobre essa versão anterior: (i) integração de múltiplas APIs de busca (Semantic Scholar, DOAJ e PubMed) em vez de depender de um único site, reduzindo a limitação de diversidade de uma fonte só; (ii) paralelização via \emph{multithreading} (\texttt{ThreadPoolExecutor}), processando múltiplos sites de busca simultaneamente; e (iii) uso de \emph{embeddings} de sentença e similaridade de cosseno para ranquear a relevância semântica dos artigos recuperados, combinado a uma etapa de scraping de reforço para localizar PDFs não diretamente acessíveis e a uma rotina de deduplicação de metadados. A arquitetura segue um modelo cliente-servidor, com interface em Streamlit e backend em Python. Os autores reportam um aumento de 64,38\% no número de artigos de referência relevantes recuperados e uma redução de 59,78\% no tempo médio de execução em relação ao sistema anterior (em números absolutos, uma média de 67 PDFs relevantes e acessíveis recuperados por consulta, contra apenas 5 no sistema anterior, com tempo médio de resposta de 35,03 segundos, contra 87,10 segundos). Assim como a pipeline proposta, o sistema combina múltiplas fontes heterogêneas e paraleliza a coleta; diferencia-se por usar o modelo apenas para medir a similaridade entre textos e ordenar os resultados por relevância, e não para interpretar o conteúdo dos documentos e dele extrair dados estruturados.

  \subsubsection{Simplificação Automática de Editais com LLM}
  Holanda \emph{et al.} [2026], em trabalho publicado no SBSI, avaliam uma pipeline baseada em LLM que gera, em modo \emph{zero-shot}, versões em linguagem simples de editais e comunicados administrativos de uma universidade pública, hoje reescritos manualmente pela instituição, com custo de tempo e retrabalho. A etapa de entrada coincide com a do sistema proposto: os editais publicados on-line são extraídos e convertidos de PDF para texto antes de irem para o \emph{prompt}. Os autores usam um \emph{prompt} fixo, fragmentam o texto em blocos de até 4.000 caracteres e fixam a temperatura em zero para garantir replicabilidade. A avaliação combina índices estatísticos de legibilidade e métricas morfossintáticas do projeto NILC-Metrix com um questionário aplicado por conveniência a docentes e estudantes que leem ou submetem editais ao menos uma vez por ano. Dois resultados são especialmente relevantes para este trabalho: na comparação entre modelos, o Gemini Flash superou o Gemini Pro em todas as métricas de legibilidade, e entre os modelos abertos o Qwen2.5:14b foi o que mais se aproximou; e vários modelos alucinaram durante a geração (produzindo textos ininteligíveis, repetindo o \emph{prompt} ou respondendo em outro idioma), o que exigiu revisão manual de todas as saídas e o registro do número de documentos alucinados por modelo. O estudo reforça, para o contexto brasileiro de documentos administrativos públicos, tanto a escolha de um modelo da linha ``Flash'' quanto a necessidade de uma camada de validação sobre a saída do LLM, ambas presentes no sistema proposto.

  \subsubsection{Agente de IA para Validação de Documentos Fiscais}
  Utyiama \emph{et al.} [2026], também no SBSI, descrevem o Smart OCR, um agente de IA baseado em LLMs que extrai campos-chave de documentos fiscais e financeiros em PDF (nota fiscal eletrônica, boleto, DANFE, DARF) e valida cada documento contra um documento de referência (\emph{ground truth}). O problema de partida é o mesmo do sistema proposto: a conferência manual desses documentos consome tempo e está sujeita a erros, e a diversidade de formatos e \emph{layouts} torna as abordagens por regras fixas pouco flexíveis. A arquitetura é orientada a filas: um lote de documentos entra em uma fila, agentes de extração aplicam OCR e transformam o conteúdo em JSON, e agentes de revisão comparam esse JSON com o documento de referência. O \emph{prompt} de extração segue princípios de clareza, contexto e estrutura, com instruções explícitas sobre o tipo de documento e os campos a extrair, restrição de formato de saída em JSON, instruções de \emph{fallback} (``caso o campo não seja encontrado, retorne null''), \emph{few-shot prompting} e verificações de consistência entre campos; a saída é ainda validada por um objeto de transferência de dados que atua como \emph{guardrail}. Os autores relatam como desafios recorrentes a perda de precisão do OCR em documentos digitalizados de baixa qualidade e os erros semânticos ocasionais do LLM quando o contexto é ambíguo, além da necessidade de um \emph{ground truth} confiável. A escolha dos modelos (Claude 3.7 Sonnet na extração, pela relação entre custo e desempenho, e Claude 3.5 Sonnet na revisão, pela estabilidade) e a preocupação com privacidade dos dados (APIs corporativas que vetam o uso dos dados para retreinamento, em conformidade com a LGPD) espelham critérios adotados no sistema proposto. O sistema foi avaliado em um estudo de caso exploratório com cinco analistas e coordenadores de um instituto de P\&D, ao longo de duas semanas, por meio de um questionário em escala Likert de cinco pontos: a utilidade percebida (M=4,73) e a intenção de recomendação foram altas, enquanto a usabilidade (M=3,90) e a satisfação inicial (M=3,80) ficaram moderadas; os autores tratam o resultado como generalização analítica para contextos organizacionais semelhantes, e não como inferência estatística.

  \subsubsection{Outro Trabalho Consultado}
  Xu \emph{et al.} [2024] apresentam uma revisão sobre o uso de LLMs em contextos educacionais, consultada aqui apenas como referência de enquadramento teórico geral. Dela se aproveitam, na Seção~\ref{sec:conceitos-fundamentais}, a caracterização de LLM como modelo baseado em processamento de linguagem natural e treinado por aprendizado auto-supervisionado em grande corpus, o crescimento acentuado da escala desses modelos e o paradigma de pré-treinamento seguido de ajuste fino.

  \subsection{Comparação dos Trabalhos Relacionados}\label{sec:comparacao-dos-trabalhos-relacionados}
  Nenhum dos trabalhos consultados trata especificamente de editais de fomento público em português, nem combina, em uma única pipeline de produção, os quatro elementos centrais do sistema proposto: (i) coleta contínua e automatizada de múltiplas fontes governamentais heterogêneas; (ii) extração de texto com \emph{fallback} de OCR para documentos sem camada de texto nativa; (iii) estruturação via LLM com vocabulário controlado; e (iv) notificação proativa sensível a prazos. Krc e Schlaack [2023] compartilham o desafio de escala e heterogeneidade de fontes governamentais, mas não empregam um LLM em nenhuma etapa, apoiando-se em conferência humana. Moreira-Filho \emph{et al.} [2025] e Chen \emph{et al.} [2025] empregam LLMs para extração estruturada e reportam métricas quantitativas de acurácia sobre conjuntos de teste, rigor de avaliação que o sistema proposto ainda não possui de forma formal (ver Seção~\ref{sec:como-a-plataforma-foi-avaliada} e Seção~\ref{sec:limitacoes}); isso configura uma lacuna concreta e um direcionamento para trabalhos futuros (Seção~\ref{sec:trabalhos-futuros}). Blázquez-Ochando, Prieto-Gutiérrez e Ovalle-Perandones [2025] usam o LLM para gerar o código de coleta, e não para interpretar o conteúdo coletado, e Fahrudin \emph{et al.} [2025] usam o LLM apenas de forma indireta, por meio de \emph{embeddings} de similaridade para ranqueamento de relevância. Os dois trabalhos do SBSI aproximam-se pelo domínio: Holanda \emph{et al.} [2026] aplicam um LLM a editais de uma universidade pública brasileira, ainda que para simplificação e não para extração estruturada, e Utyiama \emph{et al.} [2026] adotam uma arquitetura de extração e validação de campos de documentos em PDF próxima da do sistema proposto, porém no domínio fiscal e a partir de documentos já recebidos pelo fluxo de trabalho da organização, sem coleta automatizada das fontes.

  \section{Metodologia}\label{sec:metodologia}

  \subsection{Área de Estudo / Contexto de Aplicação}
  O sistema proposto foi desenvolvido no contexto de um estágio, com o objetivo de automatizar o monitoramento de editais de fomento público relevantes para a instituição. O sistema monitora quatro fontes institucionais brasileiras: FACEPE, CNPq, CAPES e FINEP.

  \subsection{Base de Dados}\label{sec:base-de-dados}
  As ``bases de dados'' deste trabalho são os quatro portais públicos monitorados pela pipeline, cada um com particularidades estruturais que exigiram estratégias de coleta distintas.

  \subsubsection{FACEPE}
  A coleta é feita via \texttt{requests} combinado a um parser HTML próprio (BeautifulSoup), que filtra apenas os editais principais do ano-alvo configurado e os ordena por data de publicação (mais recentes primeiro).

  \subsubsection{CNPq}
  A coleta segue um fluxo em duas etapas: primeiro identifica, na página de listagem, os links que já são PDFs diretos e os que são páginas de detalhe; para estas últimas, realiza uma segunda requisição HTTP (tolerante a falhas) para extrair os links de PDF contidos na página de detalhe. Os resultados são ordenados pela data de início de inscrição, mais recentes primeiro.

  \subsubsection{CAPES}
  A coleta trata a página de índice de editais como um agregador de páginas de programa: primeiro coleta as URLs das páginas de cada programa, depois visita cada uma individualmente para extrair os documentos, filtrando pelo ano-alvo. Implementa um cliente HTTP com novas tentativas e variação de host (\texttt{www.gov.br} ou \texttt{gov.br}) para lidar com instabilidades de rede observadas nesta fonte especificamente (ver Seção~\ref{sec:limitacoes}).

  \subsubsection{FINEP}
  Diferente das demais fontes, a coleta na FINEP não envolve scraping de HTML: é feita inteiramente via API REST autenticada por OAuth2 (fluxo \emph{client credentials}), consultando chamadas com situação ``aberta'' e, para cada uma, buscando o documento cujo nome contenha ``edital'' (com fallback para o primeiro documento em formato PDF disponível).

  \subsection{Desenvolvimento do Sistema}

  \subsubsection{Levantamento de Requisitos}
  A solução precisava atender a requisitos de extensibilidade (permitir adicionar novas fontes sem alterar o núcleo do sistema), resiliência a falhas de rede e de limite de uso de API externa (rate limiting do Gemini), idempotência (evitar reprocessar ou notificar duplicadamente o mesmo edital) e operação desassistida (execução diária automatizada, sem intervenção manual).

  \subsubsection{Modelagem da Solução}
  Esses requisitos motivaram um padrão de registro de fontes (\texttt{SOURCE\_REGISTRY}), no qual cada fonte expõe uma interface comum (chave, rótulo, URL base e as funções de coleta \texttt{collect\_calls}/\texttt{collect\_links}), permitindo que a camada de orquestração itere sobre as fontes de forma genérica. O README do projeto documenta esse contrato de cinco \emph{exports} obrigatórios para a adição de uma nova fonte, evidenciando a preocupação de extensibilidade desde a concepção.

  \subsubsection{Implementação}
  A implementação seguiu a divisão em camadas detalhada na Seção~\ref{sec:arquitetura-do-sistema}, com a orquestração (\texttt{pipeline\_runner.py}) coordenando sequencialmente coleta, extração, análise, persistência e notificação para cada fonte selecionada; a Figura~\ref{fig:pipeline} apresenta a visão geral desse fluxo, detalhado etapa a etapa na Seção~\ref{sec:pipeline-de-processamento}.

  \subsubsection{Tecnologias e Ferramentas Utilizadas}
  A pipeline foi implementada em \textbf{Python}, utilizando \texttt{requests} e \texttt{BeautifulSoup}/\texttt{lxml} para coleta via web scraping (FACEPE, CNPq, CAPES), um cliente OAuth2 próprio para a API REST da FINEP, \texttt{pdfplumber} para extração de texto de PDF, \texttt{pytesseract} e \texttt{pypdfium2} como fallback de OCR, o SDK \texttt{google-genai} para integração com o modelo de linguagem, \texttt{pymongo} para persistência em MongoDB e \texttt{smtplib} para notificação por e-mail. A interface web foi implementada em \textbf{Next.js 14} (App Router) com \textbf{React 18} e \textbf{TypeScript}, consumindo a mesma base MongoDB da pipeline.

  A escolha do \textbf{Gemini 2.5-Flash} como modelo de linguagem principal foi motivada por quatro fatores. Primeiro, o equilíbrio entre custo e desempenho: modelos mais robustos, como o GPT-4, apresentam custo por token elevado para uma aplicação de processamento recorrente e em escala, enquanto modelos open-source, como o Llama-3, deslocam o custo para infraestrutura de inferência dedicada (GPUs, gerenciamento de servidores), que não se justificava para o volume da aplicação. Segundo, a janela de contexto de até 1 milhão de tokens permite que um edital seja analisado integralmente, sem a necessidade de fragmentação (\emph{chunking}) do documento em blocos menores, técnica que, quando necessária, pode ocasionar perda de contexto semântico e inconsistências entre trechos processados separadamente. Terceiro, a baixa latência de inferência (``Flash''), relevante para um pipeline que processa múltiplos documentos em sequência. Quarto, a integração simplificada por meio do SDK \texttt{google.genai}, com documentação e suporte oficial consolidados.

  Ainda assim, a arquitetura foi desenhada seguindo o princípio de desacoplamento entre a lógica de aplicação e o modelo de linguagem utilizado: a função de chamada ao LLM (\texttt{call\_gemini}, em \texttt{pipeline/pdf\_pipeline/analyzer.py}) é isolada do restante do fluxo de orquestração e o identificador do modelo fica em uma única constante (com \emph{fallback} automático para uma versão equivalente da linha ``Flash'' caso a versão fixada fique indisponível), de modo que a substituição do modelo ou do provedor exija alterações localizadas, sem refatoração do restante do sistema.

  \subsection{Arquitetura do Sistema}\label{sec:arquitetura-do-sistema}

  \subsubsection{Camada de Apresentação}\label{sec:camada-de-apresentacao}
  Implementada como uma aplicação Next.js (App Router) com React e TypeScript, consumindo a mesma base MongoDB da pipeline por meio de rotas de API (\emph{Route Handlers}) e de componentes de servidor (\emph{Server Components}). A identificação do pesquisador é feita sem autenticação: basta informar um e-mail, guardado no navegador (\texttt{localStorage}), que passa a ser associado à marcação de interesse (alternância de ``estrela'') em editais específicos, restringindo o envio de lembretes de prazo apenas a esses editais marcados; a marcação em si é aplicada de forma otimista na interface (a resposta visual é imediata) e persistida de forma assíncrona no banco, com reversão automática em caso de falha de gravação.

  A listagem de editais é paginada por cursor (o \texttt{id} do último item da página anterior), 20 registros por vez, com carregamento incremental por rolagem infinita (\emph{scroll infinito}) por meio da API \texttt{IntersectionObserver} do navegador; contadores agregados (total de editais monitorados, quantidade marcada como de interesse e quantidade com prazo iminente) são calculados a partir da base completa no servidor, e não apenas dos itens já carregados na página, de modo a permanecerem corretos independentemente de quanto da lista o pesquisador já rolou. Sobre essa listagem paginada, a interface oferece dois filtros combináveis: por fonte (seleção múltipla entre FACEPE, CNPq, FINEP e CAPES) e por texto (busca pelo título do edital), aplicados em conjunto com a alternância entre ``todos os editais'' e ``apenas os marcados como interesse''.

  Cada edital, no card da listagem, expõe um link direto para o documento oficial em PDF e dá acesso a uma página de detalhamento dedicada (rota \texttt{/editais/[id]}), que reapresenta as mesmas informações estruturadas exibidas na notificação por e-mail (público-alvo, resumo, critérios de elegibilidade, cronograma e observações) em uma página web nativa em React, e não como uma reutilização direta do HTML gerado pela pipeline Python para o e-mail, já que este último depende de recursos exclusivos de clientes de e-mail (imagens anexadas via \texttt{cid:}, layout em tabelas) incompatíveis com uma página web comum.

  Além da marcação de interesse por edital, a interface oferece uma tela de preferências (``Minhas áreas de interesse''), na qual cada docente seleciona, a partir do mesmo vocabulário controlado usado pela extração (Seção~\ref{sec:transformacao}), as áreas de interesse e os segmentos que deseja acompanhar. Essas preferências são gravadas na coleção \texttt{docentes} do MongoDB e são a base do direcionamento automático de editais descrito na Seção~\ref{sec:direcionamento-automatico-de-editais}.

  \subsubsection{Camada de Negócio}
  Concentra a orquestração da pipeline, no diretório \texttt{pipeline/orchestration/}. O \texttt{pipeline\_runner.py} coordena o fluxo completo de uma fonte (coleta, extração, análise, persistência e notificação); \emph{runners} dedicados tratam os lembretes de prazo (\texttt{deadline\_reminder\_runner.py}) e a varredura avulsa de notificação por área e segmento (\texttt{new\_edital\_notify\_runner.py}). O casamento entre um edital e os docentes aderentes fica isolado em dois módulos, um que decide \emph{quem} deve ser avisado (\texttt{docente\_match.py}) e outro que decide \emph{como} enviar, com as travas de volume (\texttt{docente\_notification\_service.py}). Módulos auxiliares cuidam do registro de fontes (\texttt{source\_registry.py}), das configurações de execução (\texttt{settings.py}), da política de novas tentativas (\texttt{retry\_policy.py}) e da conversão da data limite de submissão (\texttt{date\_parser.py}). É nesta camada que residem a política de novas tentativas com espera progressiva diante de erros HTTP 429 (limite de taxa) e 503 (indisponibilidade) do provedor do LLM e o disjuntor de quota (\emph{circuit breaker}) que interrompe o processamento das fontes restantes caso a cota diária do Gemini se esgote, evitando falhas em cascata no restante da execução.

  \subsubsection{Camada de Processamento de PDF}
  Reúne, em \texttt{pipeline/pdf\_pipeline/}, o módulo de extração de texto (\texttt{extractor.py}) e o de análise via LLM (\texttt{analyzer.py}), ambos detalhados na Seção~\ref{sec:pipeline-de-processamento}.

  \subsubsection{Coletores por Fonte}
  Cada fonte é um pacote independente em \texttt{pipeline/sources/} (\texttt{facepe/}, \texttt{cnpq/}, \texttt{finep/}, \texttt{capes/}), com separação interna entre constantes, cliente HTTP ou de API, \emph{parser}, regras de negócio da fonte e orquestração da coleta. Cada pacote expõe, pelo seu \texttt{\_\_init\_\_.py}, a mesma interface pública (chave, rótulo, URL base e as funções de coleta), consumida genericamente pelo \texttt{source\_registry.py}. A função efetivamente usada pela orquestração é a \texttt{collect\_calls}, que devolve pares \texttt{(url\_pdf, data\_publicacao)} na ordem de prioridade da fonte; a \texttt{collect\_links} é um invólucro que retorna apenas as URLs. Adicionar uma fonte se resume a criar o pacote e registrá-lo. As estratégias de coleta e priorização de cada fonte são as descritas na Seção~\ref{sec:base-de-dados}.

  \subsubsection{Camada de Dados}\label{sec:camada-de-dados}
  Implementada sobre MongoDB, no diretório \texttt{pipeline/db/}. Todos os editais, de qualquer fonte, ficam em uma única coleção \texttt{editais}, na qual cada documento se identifica pelo campo \texttt{fonte} (\texttt{facepe}, \texttt{cnpq}, \texttt{finep} ou \texttt{capes}); até julho de 2026 havia uma coleção por fonte, unificadas depois por um script de migração de uso único. Os índices são um sobre \texttt{url\_pdf} (único), que serve simultaneamente como chave de deduplicação e como elo de integração entre a pipeline Python e a interface Next.js, e mais dois sobre \texttt{data\_limit\_submissao} e \texttt{fonte}, usados pelas consultas de lembrete e de cobertura. Uma segunda coleção, \texttt{docentes}, guarda o e-mail e as preferências (áreas de interesse e segmentos) de cada docente; ela é escrita pela interface e por um script de carga que importa a base a partir de um formulário de experiência profissional de docentes e servidores da UPE, e é apenas lida pela pipeline. As informações de controle de envio (\texttt{interessados}, \texttt{deadline\_reminder} e \texttt{match\_notification}) são gravadas como subdocumentos no próprio registro do edital, mantendo pipeline e interface desacopladas e comunicando-se apenas pela base compartilhada.

  \subsubsection{Direcionamento Automático de Editais}\label{sec:direcionamento-automatico-de-editais}
  O mecanismo de correspondência entre um edital e os pesquisadores aderentes, previsto pela Hipótese 1 (Seção~\ref{sec:introducao}), está implementado como uma etapa da própria pipeline e corresponde ao bloco inferior da Figura~\ref{fig:pipeline}. Ao salvar um edital novo com análise sem erro e prazo em aberto, o \texttt{pipeline\_runner.py} cruza as áreas de interesse e os segmentos que a LLM classificou naquele edital com as preferências cadastradas por cada docente (Seção~\ref{sec:camada-de-apresentacao}). A regra de correspondência é disjuntiva: basta uma área \emph{ou} um segmento em comum para o docente entrar na lista, e um edital sem nenhuma classificação não notifica ninguém. A comparação é feita sobre valores normalizados (sem espaços e sem diferença de caixa), para tolerar divergências de digitação entre o cadastro e o vocabulário controlado. O envio é feito em cópia oculta, uma mensagem por edital com todos os docentes que casam endereçados em \emph{blind carbon copy} (BCC). Sobre a lista de correspondências incidem travas de volume (no máximo três notificações para o mesmo docente e trezentos destinatários no total por execução) e um controle de idempotência gravado no próprio documento do edital (\texttt{match\_notification.sent\_emails}), de modo que reprocessar uma fonte não gere e-mail repetido; uma varredura avulsa (\texttt{main.py --notify-editais}), também agendada diariamente via GitHub Actions, cobre editais que já existiam quando o docente cadastrou as preferências ou cujo envio foi interrompido pelas travas de volume. O que ainda não foi realizado é a avaliação quase-experimental dos efeitos desse direcionamento sobre recursos captados e visibilidade institucional (Seção~\ref{sec:limitacoes}).

  \subsection{Pipeline de Processamento}\label{sec:pipeline-de-processamento}

  A Figura~\ref{fig:pipeline} apresenta uma visão geral do fluxo completo do sistema proposto, desde a publicação de um edital em um portal institucional até a notificação dos interessados, e organiza o processo em dois blocos encadeados. O primeiro bloco, de monitoramento e extração, parte da coleta das fontes (Seção~\ref{sec:base-de-dados}), baixa e pré-processa cada documento em PDF com extração de texto e \emph{fallback} de OCR (Seção~\ref{sec:extracao}), submete o texto ao LLM por meio de um \emph{prompt} estruturado (Seção~\ref{sec:transformacao}), aplica uma camada defensiva de validação e normalização sobre os campos retornados (Seção~\ref{sec:transformacao}) e persiste o resultado como JSON no MongoDB (Seções~\ref{sec:camada-de-dados} e~\ref{sec:visualizacao-saida}); falhas de interpretação da resposta do modelo ou inconsistências detectadas na validação realimentam etapas anteriores, reprocessando o documento (Seção~\ref{sec:processamento}). O segundo bloco, de notificações, é disparado a cada edital novo persistido: cruza os campos de classificação extraídos com as preferências cadastradas pelos docentes segundo a regra de correspondência descrita na Seção~\ref{sec:direcionamento-automatico-de-editais}, monta o alerta com o resumo estruturado do edital, envia as mensagens por e-mail nos três fluxos previstos (Seção~\ref{sec:visualizacao-saida}) e registra cada envio no próprio documento do edital para garantir idempotência (Seção~\ref{sec:camada-de-dados}). As subseções seguintes detalham as etapas do primeiro bloco.

  \begin{figure}[t]
  \centering
  \includegraphics[width=\textwidth]{figures/Figura 1.png}
  \caption{Visão geral da \emph{pipeline} proposta. O bloco superior cobre o monitoramento e a extração estruturada dos editais (Seções~\ref{sec:base-de-dados} a~\ref{sec:pipeline-de-processamento}), com realimentação para reprocessamento em caso de inconsistências; o bloco inferior cobre o direcionamento automático e a notificação dos interessados (Seções~\ref{sec:direcionamento-automatico-de-editais} e~\ref{sec:visualizacao-saida}).}
  \label{fig:pipeline}
  \end{figure}

  \subsubsection{Extração}\label{sec:extracao}
  O documento em PDF é baixado em memória e processado com \texttt{pdfplumber}. Quando essa extração retorna um documento sem nenhum caractere reconhecível (cenário observado tanto em digitalizações quanto em PDFs ``impressos'' a partir de visualizadores que bloqueiam a seleção de texto, nos quais cada letra é armazenada como forma vetorial e não como glifo de fonte), o sistema recorre a um fallback de OCR: cada página é renderizada como imagem (via \texttt{pypdfium2}) e processada com \texttt{pytesseract} (idioma português), limitado às primeiras 20 páginas do documento por padrão, dado o custo computacional mais alto do OCR em relação à extração direta.

  \subsubsection{Transformação}\label{sec:transformacao}
  O texto extraído é inserido em um prompt estruturado que instrui o LLM a atuar como um ``analista de editais'' e retornar exclusivamente um objeto JSON, sem inventar informações não presentes no texto. O objeto reúne o título oficial, o público-alvo, uma descrição, duas listas de critérios de elegibilidade (uma para o público-alvo e outra para o proponente), uma lista de observações, o cronograma e a data limite de submissão, além dos dois campos de classificação usados no direcionamento. Esses dois campos, \texttt{areas\_interesse} (oito valores possíveis) e \texttt{segmentos} (dez valores possíveis), são restritos a vocabulários controlados fornecidos no próprio prompt, que também exige pelo menos um valor em cada e classificação baseada em evidência explícita no edital; o campo de prazo deve seguir o formato ISO (\texttt{AAAA-MM-DD}). Após a resposta do modelo, o sistema aplica uma camada defensiva de validação: preenche campos ausentes, remove dos campos de classificação os valores fora do vocabulário controlado e elimina duplicatas em \texttt{areas\_interesse}, \texttt{segmentos} e \texttt{cronograma}.

  \subsubsection{Processamento}\label{sec:processamento}
  A chamada ao modelo (Gemini 2.5-Flash, via SDK \texttt{google.genai}, com saída forçada em JSON) é envolvida por uma política de novas tentativas: em caso de erro 429 com tempo de espera sugerido pela própria API, o sistema aguarda o tempo indicado; em caso de esgotamento genuíno de cota diária, desiste imediatamente (sem novas tentativas); em caso de erro 503, aplica espera progressiva. Falhas de interpretação do JSON retornado disparam uma segunda tentativa de extração via expressão regular antes de registrar o documento como erro de análise.

  \subsubsection{Visualização / Saída}\label{sec:visualizacao-saida}
  Os dados estruturados são persistidos no MongoDB (com controle de duplicidade por \texttt{url\_pdf}) e tornam-se imediatamente consultáveis pela interface web. Quando um edital novo é salvo, ou um edital antes registrado como erro é reprocessado com sucesso, e o prazo de submissão ainda não expirou, o sistema dispara notificações por e-mail, todas com cabeçalho de identidade visual institucional (logos embutidos via \texttt{cid:}, para compatibilidade com clientes que bloqueiam imagens externas). São três fluxos distintos: um aviso institucional, enviado a um endereço que recebe todos os editais; o direcionamento por área e segmento (Seção~\ref{sec:direcionamento-automatico-de-editais}), enviado em cópia oculta (BCC) a cada docente cujas preferências casam com o edital; e, mais adiante no ciclo, os lembretes de prazo nos marcos configurados (por padrão trinta, quinze e sete dias), enviados individualmente a cada pesquisador que marcou aquele edital específico como de interesse.

  \subsection{Como a Plataforma foi Avaliada}\label{sec:como-a-plataforma-foi-avaliada}

  \subsubsection{Validação Incidental em Produção}\label{sec:validacao-incidental-em-producao}
  A validação da solução, até o momento, foi conduzida de forma incidental, isto é, orientada pela observação do comportamento em produção, e não por um protocolo formal de avaliação com métricas de acurácia sobre um conjunto de teste rotulado (diferentemente da abordagem adotada, por exemplo, por Moreira-Filho \emph{et al.} [2025] e por Chen \emph{et al.} [2025], discutidos na Seção~\ref{sec:trabalhos-relacionados}). Essa validação incidental está documentada no \texttt{TROUBLESHOOTING.md} do projeto e inclui, entre outros casos: a identificação de que documentos sem camada de texto extraível ficavam permanentemente presos com status de erro, o que motivou a adição do fallback de OCR (Seção~\ref{sec:extracao}); e a identificação de que editais reprocessados com sucesso (de erro para status ``ok'') não disparavam notificação por e-mail, por uma condição que verificava apenas o status ``inserido'' e não ``atualizado''. Esse segundo caso foi corrigido após observação do comportamento real do sistema. A ausência de métricas formais de acurácia, discutida como limitação na Seção~\ref{sec:limitacoes}, motivou o desenho dos dois protocolos formais descritos a seguir, um para cada hipótese do projeto.

  \subsubsection{Protocolo de Avaliação da Confiabilidade da Extração (Hipótese 2)}\label{sec:protocolo-de-avaliacao-da-confiabilidade-da-extracao-hipotese-2}
  A Hipótese 2 (Seção~\ref{sec:introducao}) afirma que a extração dos campos essenciais de um edital, realizada por um LLM orientado por prompt estruturado com vocabulário controlado (Seção~\ref{sec:transformacao}), atinge precisão e recall suficientes para sustentar o direcionamento automático investigado pela Hipótese 1. Ela é avaliada por um protocolo no qual um avaliador humano lê o documento PDF original de uma amostra de editais de cada fonte lado a lado com a mesma informação exibida pela interface web (Seção~\ref{sec:camada-de-apresentacao}): (i) órgão e prazo final de submissão; (ii) público-alvo; (iii) resumo do edital; (iv) áreas de interesse; (v) segmentos; (vi) critérios do público-alvo; (vii) quem pode submeter; (viii) cronograma; e (ix) observações. Para cada um desses nove campos, o avaliador registra três julgamentos padronizados: (a) uma nota de fidelidade ao documento original, em escala de 1 a 5; (b) a ocorrência de alucinação, isto é, informação presente na saída da IA sem correspondência no documento original; e (c) a ocorrência de omissão, isto é, informação relevante presente no documento e não capturada pela extração. O instrumento de coleta (Seção~\ref{sec:instrumentos-de-coleta}) é preenchido um formulário por edital avaliado; dele serão derivadas, para a versão final deste artigo, métricas de precisão, recall e F1 por campo, seguindo a mesma orientação de avaliação quantitativa adotada por Moreira-Filho \emph{et al.} [2025] e por Chen \emph{et al.} [2025] (Seção~\ref{sec:trabalhos-relacionados}).

  \subsubsection{Protocolo de Avaliação do Direcionamento Automático (Hipótese 1)}\label{sec:protocolo-de-avaliacao-do-direcionamento-automatico-hipotese-1}
  A Hipótese 1 (Seção~\ref{sec:introducao}) afirma que o direcionamento automático de cada edital aos pesquisadores, departamentos ou grupos de pesquisa mais aderentes de sua área aumenta o tempo de preparação disponível e está associado a mais recursos de fomento captados e maior visibilidade institucional. O mecanismo de correspondência (\emph{matching}) entre os campos de classificação extraídos pela pipeline e o cadastro de preferências de docentes já está implementado e integrado à execução (Seção~\ref{sec:direcionamento-automatico-de-editais}): cada edital novo é direcionado aos docentes cujas áreas ou segmentos coincidem com os dele. O que ainda não foi conduzido é o experimento controlado que a hipótese exige, comparando o comportamento de submissão e os resultados de captação de grupos com e sem esse direcionamento ao longo de um ou mais ciclos de fomento (Seção~\ref{sec:trabalhos-futuros}). Como etapa preliminar a essa comparação, foi desenhado um instrumento de percepção, aplicado a docentes, coordenadores de grupo de pesquisa e gestores. O instrumento mede, em escala Likert de 5 pontos, o quanto o direcionamento automático é percebido como capaz de aumentar o tempo disponível para organização de propostas, a competitividade das submissões, o volume de recursos captados e a visibilidade institucional, além de levantar uma linha de base do acompanhamento manual atual (frequência de perda de prazos e tempo semanal dedicado à busca de editais).

  \subsubsection{Instrumentos de Coleta}\label{sec:instrumentos-de-coleta}
  Os dois protocolos acima foram operacionalizados como formulários eletrônicos estruturados (Google Forms), gerados programaticamente por um script (Google Apps Script) versionado junto ao projeto, de modo a garantir que as perguntas sobre áreas de interesse e segmentos usem exatamente o mesmo vocabulário controlado empregado pela extração (Seção~\ref{sec:transformacao}), evitando divergência entre os termos avaliados e os termos efetivamente produzidos pelo sistema. \textcolor{red}{[ATUALIZAR: inserir aqui, na versão final, o número de respondentes/editais avaliados por instrumento e o período de coleta, assim que a avaliação for aplicada.]}

  \section{Resultados}\label{sec:resultados}

  \subsection{Implementação da Solução}
  A pipeline é executada por meio de uma interface de linha de comando (\texttt{pipeline/main.py}), que aceita a fonte (\texttt{--source}) e um limite de editais a processar por execução (\texttt{--limit}), além de modos dedicados para o envio de lembretes de prazo (\texttt{--run-reminders}, com marcos configuráveis, por padrão 30, 15 e 7 dias) e para a varredura avulsa de direcionamento por área e segmento (\texttt{--notify-editais}). Em produção, a execução é automatizada via GitHub Actions por três rotinas agendadas: a pipeline principal (\texttt{pipeline\_runner\_daily.yml}), que itera sobre as quatro fontes com um limite conservador de processamento por fonte (para preservar a cota diária do Gemini); os lembretes de prazo (\texttt{deadline-reminders.yml}), uma hora depois; e o \emph{backfill} do direcionamento por área e segmento (\texttt{notify-editais-daily.yml}), mais uma hora depois, que recupera os casamentos ainda não notificados. Cada rotina inclui a liberação temporária do IP do executor no allowlist do MongoDB Atlas (removida ao final da execução); a pipeline principal inclui ainda um disjuntor de quota em nível de shell, que interrompe o processamento das fontes restantes caso a cota do Gemini já tenha se esgotado em uma fonte anterior na mesma execução.

  \subsection{Interface do Sistema}
  O sistema proposto disponibiliza uma interface web pública, hospedada em \texttt{editais.upe.br} e identificada como ``Plataforma de Monitoramento de Editais'', que dá ao pesquisador acesso consultável a tudo que a pipeline coletou e estruturou, sem exigir autenticação. Um selo no topo indica se os dados exibidos são reais (conexão ativa ao MongoDB, rótulo ``ao vivo'') ou de demonstração (ausência de configuração de banco). As subseções seguintes percorrem as principais telas, referenciadas nas Figuras~\ref{fig:ui-preferencias} a \ref{fig:ui-detalhe}.

  \subsubsection{Identificação e Preferências de Acompanhamento}
  A identificação do pesquisador é feita apenas por endereço de e-mail, sem senha ou cadastro; o e-mail informado fica guardado no navegador e passa a ser exibido no cabeçalho (``marcando interesse como \ldots'', com a opção de trocar). A partir dele, o painel ``Minhas áreas de interesse'' permite selecionar, a partir exatamente do mesmo vocabulário controlado usado pela extração (Seção~\ref{sec:transformacao}), as oito áreas de interesse e os dez segmentos que o pesquisador deseja acompanhar. Ao salvar, essas preferências são gravadas na coleção \texttt{docentes} e passam a alimentar o direcionamento automático descrito na Seção~\ref{sec:direcionamento-automatico-de-editais}: sempre que a pipeline classificar um edital novo com pelo menos uma área ou um segmento em comum com o perfil, o pesquisador recebe uma notificação por e-mail. A Figura~\ref{fig:ui-preferencias} mostra essa tela, com o cabeçalho de identificação e o painel de preferências.

  \begin{figure}[t]
  \centering
  \includegraphics[width=\textwidth]{figures/Figura 2.png}
  \caption{Tela inicial da interface web (\texttt{editais.upe.br}): identificação por e-mail (sem senha) e painel ``Minhas áreas de interesse'', com a seleção das áreas e dos segmentos do vocabulário controlado que definem o direcionamento automático (Seção~\ref{sec:direcionamento-automatico-de-editais}).}
  \label{fig:ui-preferencias}
  \end{figure}

  \subsubsection{Painel de Indicadores e Filtros}
  Abaixo das preferências, três indicadores agregados resumem o estado do acompanhamento: o total de editais monitorados nas quatro fontes, a quantidade marcada como de interesse pelo pesquisador e a quantidade, entre as marcadas, com prazo de submissão em até sete dias. Esses contadores são calculados no servidor a partir da base completa, e não apenas dos itens já carregados na tela. Uma barra de ferramentas oferece, ainda, a alternância entre a visão ``Todos'' e ``Meus interesses'', uma busca por texto do título e um filtro por fonte (seleção múltipla entre FACEPE, CNPq, FINEP e CAPES); os filtros são combináveis entre si e com a visão de interesse.

  \subsubsection{Listagem de Editais}
  Os editais são apresentados em uma grade de \emph{cards}, ordenados pela data de publicação na fonte de origem (mais recentes primeiro, seguindo a ordem do portal oficial; quando essa data não está disponível, usa-se a data de registro na base), com os editais de prazo já expirado removidos da vitrine. O carregamento é incremental por rolagem infinita, vinte editais por vez. Cada \emph{card} exibe a fonte e um código de referência, o título oficial extraído pela pipeline, o órgão responsável, o prazo final de submissão com a contagem de dias restantes (destacada por cor conforme a urgência) ou a marca ``sem prazo definido'', e as etiquetas de áreas de interesse e segmentos atribuídas pela classificação. Dois controles encerram o \emph{card}: o botão ``Acessar Edital'', que abre diretamente o PDF oficial na fonte, e o botão de alternância em forma de estrela (``marque para receber lembretes de prazo''), que inclui ou remove aquele edital específico da lista pessoal de acompanhamento do pesquisador. Marcar a estrela é o que faz o edital passar a gerar os lembretes por e-mail em D-30, D-15 e D-7 (Seção~\ref{sec:visualizacao-saida}), enviados apenas para quem marcou aquele edital. A marcação é aplicada de forma otimista na interface, com resposta visual imediata e reversão automática em caso de falha de gravação. A Figura~\ref{fig:ui-listagem} mostra a barra de filtros e a grade de \emph{cards}.

  \begin{figure}[t]
  \centering
  \includegraphics[width=\textwidth]{figures/Figura 3.png}
  \caption{Listagem de editais: barra com as abas ``Todos''/``Meus interesses'', busca por título e filtro por fonte, seguida da grade de \emph{cards}. Cada \emph{card} traz fonte, título, órgão, prazo com dias restantes, etiquetas de área e segmento, o acesso ao PDF oficial e a estrela que aciona os lembretes de prazo.}
  \label{fig:ui-listagem}
  \end{figure}

  \subsubsection{Página de Detalhamento do Edital}
  A partir de qualquer \emph{card}, o pesquisador pode abrir uma página de detalhamento dedicada (rota \texttt{/editais/[id]}), que reapresenta em uma página web nativa todos os campos que a pipeline estruturou para aquele edital (público-alvo, resumo, critérios de elegibilidade do público-alvo e do proponente, cronograma e observações), as mesmas informações enviadas no e-mail de notificação. Assim, o pesquisador tem acesso ao resumo estruturado do edital sem depender de localizar o e-mail correspondente nem de ler o PDF original na íntegra. A Figura~\ref{fig:ui-detalhe} mostra essa página.

  \begin{figure}[t]
  \centering
  \includegraphics[width=\textwidth]{figures/Figura 4.png}
  \caption{Página de detalhamento de um edital (rota \texttt{/editais/[id]}), com os campos estruturados pela pipeline (público-alvo, resumo, critérios de elegibilidade, cronograma e observações) apresentados em página web nativa.}
  \label{fig:ui-detalhe}
  \end{figure}

  \subsubsection{Uso em Dispositivos Móveis}
  A identificação por e-mail também é solicitada de forma contextual: ao tentar marcar interesse em um edital sem ainda ter informado o e-mail, a própria notificação de confirmação exibe o campo para isso, evitando que o pesquisador precise procurar o campo de identificação em outra parte da página. Em telas de celular, essas notificações de confirmação assumem o formato de faixa fixada na borda inferior da tela, em vez do selo flutuante usado em telas maiores, ajuste motivado por observação direta de uso da interface em dispositivo móvel.

  \subsection{Estudo de Caso}
  Cada fonte monitorada é apresentada como um estudo de caso, cruzando o conjunto de editais que o coletor prioriza no portal de origem (as chamadas do ano corrente, ou as mais recentes no caso da FINEP) com o número de documentos já processados e armazenados pela pipeline. Apresenta-se também o acervo acumulado da fonte na base, que inclui documentos de execuções e de anos anteriores.

  Os números a seguir foram levantados em 3 de setembro de 2026 com o utilitário \texttt{check\_mongo\_coverage.py}, que cruza os links coletados de cada portal com os documentos correspondentes na coleção \texttt{editais}, e com uma consulta agregada direta à mesma coleção; a Tabela~\ref{tab:cobertura} os consolida. Na data da verificação, o portal do CNPq não listava nenhuma chamada aberta, de modo que a análise ao vivo dessa fonte se limita ao que já está consolidado na base.

  \subsubsection{FACEPE}
  A coleta retorna, para o ano-alvo corrente, 25 editais principais, todos os 25 já processados com sucesso, nenhum pendente ou com erro. No acervo acumulado, a base guarda 85 documentos da FACEPE (81 com sucesso e 4 com erro de processamento).

  \subsubsection{CAPES}
  São 28 editais na seção de editais abertos do ano-alvo, todos os 28 processados com sucesso. O acervo acumulado tem 31 documentos da CAPES, todos com status de sucesso.

  \subsubsection{CNPq}
  No momento da verificação, o portal monitorado não listava nenhuma chamada aberta (a página retorna explicitamente a mensagem ``Nenhuma chamada aberta.''). O acervo guarda 1 edital da CNPq, coletado em um ciclo anterior e processado com sucesso. O volume estruturalmente baixo dessa fonte, discutido na Seção~\ref{sec:analise-dos-resultados}, decorre da política de publicação da agência e da alternância entre períodos com e sem chamadas abertas, não de falha de coleta.

  \subsubsection{FINEP}
  A API retorna as 8 chamadas abertas mais recentes, todas já processadas com sucesso. No acervo acumulado, a base guarda 16 documentos da FINEP (15 com sucesso e 1 com erro).

  \begin{table}[t]
  \caption{Cobertura da pipeline por fonte, em 3 de setembro de 2026 (\texttt{check\_mongo\_coverage.py} e consulta agregada à coleção \texttt{editais}). \emph{Portal}: editais priorizados pelo coletor na origem; \emph{Proc.}/\emph{Erro}: dentre esses, os já persistidos com \texttt{status} \texttt{ok} e \texttt{erro}; \emph{Falta}: sem documento correspondente na base; \emph{Acervo}: total acumulado da fonte na base, incluindo execuções e anos anteriores.}
  \label{tab:cobertura}
  \centering
  \begin{tabular}{lccccc}
  \hline
  \textbf{Fonte} & \textbf{Portal} & \textbf{Proc.} & \textbf{Erro} & \textbf{Falta} & \textbf{Acervo} \\
  \hline
  FACEPE & 25 & 25 & 0 & 0 & 85 \\
  CAPES  & 28 & 28 & 0 & 0 & 31 \\
  FINEP  & 8  & 8  & 0 & 0 & 16 \\
  CNPq\textsuperscript{*} & 0 & 0 & 0 & 0 & 1 \\
  \hline
  \textbf{Total} & \textbf{61} & \textbf{61} & \textbf{0} & \textbf{0} & \textbf{133} \\
  \hline
  \end{tabular}

  \vspace{2pt}
  {\footnotesize \textsuperscript{*}Na data da verificação o portal do CNPq não listava chamadas abertas; o registro no acervo provém de coleta anterior.}
  \end{table}

  \subsection{Análise dos Resultados}\label{sec:analise-dos-resultados}
  Considerado o conjunto que cada coletor prioriza, a cobertura da pipeline é hoje completa nas fontes com chamadas abertas: dos 61 editais atualmente priorizados por FACEPE (25), CAPES (28) e FINEP (8), todos os 61 já estão processados e nenhum ficou com erro; a CNPq não tinha chamadas abertas na data da verificação e seu único edital no acervo provém de coleta anterior. Na base como um todo, são 133 editais armazenados, dos quais 5 (cerca de 3,8\%) ficaram com erro de processamento, concentrados na FACEPE (4) e na FINEP (1). A coleção de preferências de docentes, que alimenta o direcionamento automático (Seção~\ref{sec:direcionamento-automatico-de-editais}), tem 264 registros, todos com pelo menos uma área ou um segmento marcado.

  A distribuição desigual observada em execuções anteriores, com a FACEPE muito à frente das demais fontes, refletia sobretudo o acúmulo histórico e o amadurecimento desigual dos coletores, e não uma limitação estrutural. Com os coletores das quatro fontes estabilizados, o volume coletado passou a acompanhar o que cada portal efetivamente publica no período monitorado; o menor volume, no CNPq, decorre da própria política de publicação da agência no portal monitorado, e não de falha de coleta.

  \section{Discussão}\label{sec:discussao}

  \subsection{Comparação com Trabalhos Relacionados}
  Como discutido na Seção~\ref{sec:comparacao-dos-trabalhos-relacionados}, o sistema proposto ocupa um espaço pouco explorado na literatura consultada: uma pipeline de produção que combina coleta multi-fonte de documentos governamentais heterogêneos (como em \cite{b3}) com extração semântica via LLM orientado por prompt estruturado (como em \cite{b4} e em \cite{b6}). Ao mesmo tempo, o sistema ainda não incorpora o rigor de avaliação formal presente nesses dois últimos trabalhos, o que limita a comparabilidade direta de sua acurácia com a reportada por eles.

  \subsection{Contribuições da Solução}
  Entre as contribuições concretas do sistema estão: (i) uma arquitetura multi-fonte com um contrato de extensão documentado, permitindo a adição de novas fontes sem alteração do núcleo de orquestração; (ii) um fallback de OCR que recupera documentos que, de outra forma, ficariam permanentemente sem processamento; (iii) um esquema de extração com vocabulário controlado, que reduz a variabilidade e facilita o consumo posterior dos dados extraídos; (iv) um sistema de lembrete de prazos sensível a marcos configuráveis, com filtro por interesse do pesquisador; e (v) uma arquitetura desacoplada do provedor de LLM, reduzindo o custo de uma eventual troca de modelo ou fornecedor.

  \subsection{Limitações}\label{sec:limitacoes}
  O sistema monitora atualmente apenas quatro fontes, todas brasileiras, sendo três de abrangência nacional (CNPq, CAPES e FINEP) e uma de âmbito estadual (FACEPE, Pernambuco): não há cobertura de fundações estaduais de amparo à pesquisa de outras unidades da federação nem de fontes internacionais. A fonte CAPES apresenta uma instabilidade de rede não resolvida (erro de ``rede inalcançável'' a partir do ambiente de execução em nuvem), documentada no \texttt{TROUBLESHOOTING.md} como um problema em aberto. O fallback de OCR está limitado, por padrão, às primeiras 20 páginas de um documento, o que pode deixar de capturar informações relevantes em editais mais extensos. A cota diária do provedor de LLM impõe um limite conservador ao volume de documentos processados por execução, retardando a cobertura completa das fontes com maior volume de publicação. O mecanismo de direcionamento automático de editais aos interessados de cada área, central para a Hipótese 1 (Seção~\ref{sec:introducao}), está implementado e em operação (Seção~\ref{sec:direcionamento-automatico-de-editais}), mas apresenta duas limitações. A primeira é de método: a correspondência é heurística, feita pelo cruzamento de valores normalizados do vocabulário controlado, o que tolera divergências de digitação mas não capta sinonímia ou proximidade semântica entre termos; e o cadastro de perfis cobre apenas docentes individuais, ainda não departamentos nem grupos de pesquisa. A segunda é de avaliação: os efeitos desse direcionamento sobre o volume de recursos captados e a visibilidade institucional ainda não foram medidos, pois o experimento quase-experimental que a Hipótese 1 exige não foi conduzido (Seção~\ref{sec:trabalhos-futuros}). Por fim, o sistema não dispõe, até o momento, de uma avaliação formal de acurácia da extração sobre um conjunto de teste rotulado, como as reportadas pelos trabalhos discutidos na Seção~\ref{sec:trabalhos-relacionados}, de modo que sua confiabilidade tem sido inferida indiretamente, a partir da ausência de incidentes reportados, e não medida diretamente sobre um conjunto de teste. Um protocolo formal para essa avaliação, e para a avaliação complementar da Hipótese 1, já foi desenhado (Seção~\ref{sec:como-a-plataforma-foi-avaliada}) mas ainda não foi executado em escala suficiente para produzir métricas consolidadas; ambas as limitações, portanto, permanecem até a conclusão dessa coleta.

  \subsection{Impactos e Aplicações}
  Para além do escopo imediato de editais de fomento, a arquitetura do sistema proposto (coleta multi-fonte desacoplada, extração resiliente a documentos sem texto nativo e estruturação via LLM com vocabulário controlado) é potencialmente aplicável a outros cenários de monitoramento institucional de documentos públicos heterogêneos, como acompanhamento de licitações, diários oficiais ou normativas regulatórias.

  \section{Conclusão}\label{sec:conclusao}

  \subsection{Síntese dos Resultados}
  Este artigo apresentou uma pipeline automatizada que integra coleta multi-fonte, extração de texto resiliente (com fallback de OCR) e estruturação de dados via LLM para o monitoramento de editais de fomento público, concebida como base técnica para o direcionamento automático desses editais aos interessados certos de cada área (Hipótese 1, Seção~\ref{sec:introducao}). O sistema está em operação, com cobertura completa nas fontes recoletadas: dos 61 editais atualmente priorizados pelos coletores de FACEPE, CAPES e FINEP, todos os 61 já foram processados e armazenados sem erro, e a base acumula 133 editais das quatro fontes.

  \subsection{Atendimento aos Objetivos}
  Os objetivos específicos de implementação dos coletores, extração com OCR, estruturação via LLM, persistência com deduplicação e notificação por e-mail (Seção~\ref{sec:introducao}) foram atendidos e estão em operação em produção. O objetivo relativo à disponibilização de uma interface para marcação de interesse também foi atendido, com a ressalva de que sua adoção efetiva por pesquisadores ainda não foi avaliada.

  \subsection{Confirmação ou Refutação das Hipóteses}

  \subsubsection{Hipótese 1 (Direcionamento Automático)}
  A Hipótese 1 (Seção~\ref{sec:introducao}) ainda não pode ser confirmada ou refutada. O mecanismo que ela investiga, o cruzamento entre os campos de classificação extraídos pela pipeline e um cadastro de preferências de docentes, está implementado e em operação (Seção~\ref{sec:direcionamento-automatico-de-editais}): cada edital novo é direcionado aos docentes cujas áreas ou segmentos coincidem com os dele. O que ainda não foi conduzido é o experimento quase-experimental que compara o comportamento de submissão e os resultados de captação de grupos com e sem esse direcionamento ao longo de um ou mais ciclos de fomento. O instrumento de percepção descrito na Seção~\ref{sec:protocolo-de-avaliacao-do-direcionamento-automatico-hipotese-1}, aplicado a docentes, coordenadores de grupo de pesquisa e gestores, é uma etapa preliminar de avaliação, mas não substitui essa comparação quase-experimental para confirmar ou refutar a hipótese em termos de recursos efetivamente captados e de visibilidade institucional. \textcolor{red}{[ATUALIZAR com o resultado da coleta de percepção assim que disponível.]}

  \subsubsection{Hipótese 2 (Confiabilidade da Extração)}
  A Hipótese 2 (Seção~\ref{sec:introducao}) também não pode ainda ser formalmente confirmada ou refutada, pois o protocolo de avaliação campo a campo descrito na Seção~\ref{sec:protocolo-de-avaliacao-da-confiabilidade-da-extracao-hipotese-2} ainda está em fase de coleta, e não foi conduzido um experimento controlado comparando o esforço e a confiabilidade do acompanhamento manual com os do acompanhamento automatizado pela pipeline. O que se pode afirmar, a partir da observação em produção (Seção~\ref{sec:validacao-incidental-em-producao}), é que a pipeline é tecnicamente capaz de extrair os campos essenciais de um edital de forma consistente na maioria dos casos observados, e que o fallback de OCR e a política de novas tentativas reduziram a incidência de documentos sem processamento. Isso constitui evidência favorável, mas indireta, e não uma validação formal da hipótese. \textcolor{red}{[ATUALIZAR com as métricas de precisão/recall/F1 por campo assim que a coleta do formulário de avaliação for concluída.]}

  \subsection{Trabalhos Futuros}\label{sec:trabalhos-futuros}
  Como direção futura prioritária, destaca-se a avaliação dos efeitos do direcionamento automático investigado pela Hipótese 1 (Seção~\ref{sec:introducao}), cujo mecanismo de correspondência (\emph{matching}) entre os campos classificados de cada edital e o cadastro de preferências de docentes já está implementado (Seção~\ref{sec:direcionamento-automatico-de-editais}): um desenho quase-experimental que compare, para uma amostra de editais, o comportamento de submissão de grupos que recebem direcionamento automático com o de grupos que dependem apenas do acompanhamento manual atual, acompanhando ao longo de um ou mais ciclos de fomento a evolução do volume de recursos captados e de indicadores de visibilidade institucional junto às agências de fomento. Complementa essa direção a expansão do cadastro de perfis para além de docentes individuais, incorporando departamentos e grupos de pesquisa categorizados pelo mesmo vocabulário controlado usado na extração (Seção~\ref{sec:transformacao}), e a evolução da correspondência hoje heurística para uma comparação sensível a sinonímia e proximidade semântica entre termos. Complementarmente, destacam-se: a condução de uma avaliação formal de acurácia da extração, seguindo a mesma orientação de avaliação quantitativa reportada por Moreira-Filho \emph{et al.} [2025] e por Chen \emph{et al.} [2025]; a investigação de motores de OCR alternativos ao \texttt{pytesseract}, dado o limite de tokens de entrada e a sensibilidade a documentos multilíngues que Chen \emph{et al.} [2025] relatam como limitações de sistemas comparáveis baseados em OCR e LLM; a adoção de deduplicação por conteúdo (\emph{checksum}), mais robusta que a deduplicação por URL usada hoje e que a deduplicação por nome de arquivo relatada por Krc e Schlaack [2023], para lidar com revisões de um mesmo edital; a investigação de causa raiz e correção da instabilidade de rede específica da fonte CAPES; e a expansão da cobertura para novas fontes de fomento, estaduais ou internacionais.

  \begin{thebibliography}{00}
  \bibitem[Blázquez-Ochando et al. 2025]{b1} Blázquez-Ochando, M., Prieto-Gutiérrez, J. J., and Ovalle-Perandones, M. A. (2025). Prompt engineering for bibliographic web-scraping. \emph{Scientometrics}, 130:3433--3453.
  \bibitem[Chen et al. 2025]{b6} Chen, L.-C. et al. (2025). Evaluation of prompt engineering on the performance of a large language model in document information extraction. \emph{Electronics}, 14(11):2145.
  \bibitem[Fahrudin et al. 2025]{b2} Fahrudin, T. M., Funabiki, N., Brata, K. C., Naing, I., Aung, S. T., Muhaimin, A., and Prasetya, D. A. (2025). An improved reference paper collection system using web scraping with three enhancements. \emph{Future Internet}, 17(5):195.
  \bibitem[Holanda et al. 2026]{b8} Holanda, J. P., Magalhães, R., Almendra, C., and do Rêgo, L. G. C. (2026). Simplifying administrative texts for plain language using LLM: a model comparative analysis. In \emph{Anais do XXII Simpósio Brasileiro de Sistemas de Informação (SBSI)}, pages 931--947, Porto Alegre. SBC.
  \bibitem[Krc and Schlaack 2023]{b3} Krc, M. and Schlaack, A. O. (2023). Pipeline or pipe dream: Building a scaled automated metadata creation and ingest workflow using web scraping tools. \emph{Code4Lib Journal}, (58).
  \bibitem[Moreira-Filho et al. 2025]{b4} Moreira-Filho, J. T. et al. (2025). Automating data extraction from scientific literature and general PDF files using large language models and KNIME: An application in toxicology. \emph{WIREs Computational Molecular Science}, 15:e70047.
  \bibitem[Sahoo et al. 2024]{b5} Sahoo, P., Singh, A. K., Saha, S., Jain, V., Mondal, S., and Chadha, A. (2024). A systematic survey of prompt engineering in large language models: Techniques and applications. \emph{arXiv preprint arXiv:2402.07927}.
  \bibitem[Utyiama et al. 2026]{b9} Utyiama, D. M. da S., Moraes, O. S. C., Guimarães, F. A. S., Ramos, P. B. S., and Marques, L. C. (2026). Smart OCR: a vertical AI agent for financial document validation, an exploratory case study. In \emph{Anais do XXII Simpósio Brasileiro de Sistemas de Informação (SBSI)}, pages 948--967, Porto Alegre. SBC.
  \bibitem[Xu et al. 2024]{b7} Xu, H., Gan, W., Qi, Z., Wu, J., and Yu, P. S. (2024). Large language models for education: A survey. \emph{arXiv preprint arXiv:2405.13001}.
  \end{thebibliography}

  \end{document}
